\textbf{Theorem (Cauchy–Schwarz).}  
For any real numbers $a_{1},\dots ,a_{n}$ and $b_{1},\dots ,b_{n}$,
\[
\left(\sum_{i=1}^{n} a_{i}b_{i}\right)^{2}
\le
\left(\sum_{i=1}^{n} a_{i}^{2}\right)
\left(\sum_{i=1}^{n} b_{i}^{2}\right).
\]

Equality holds iff there exists $\lambda\in\mathbb R$ such that $a_{i}= \lambda b_{i}$ for every $i$ (including the trivial case $a_{i}=b_{i}=0$ for all $i$).

\begin{proof}
\begin{enumerate}
\item \textbf{Quadratic function.}  
Define
\[
f(t)=\sum_{i=1}^{n}(a_{i}-t b_{i})^{2},\qquad t\in\mathbb R .
\]
Each summand is a square, hence $f(t)\ge 0$ for all real $t$.

\item \textbf{Expansion.}  
For a single index $i$,
\[
(a_{i}-t b_{i})^{2}=a_{i}^{2}-2a_{i}b_{i}t+b_{i}^{2}t^{2}.
\]
Summing over $i$ gives
\begin{align*}
f(t)&=\sum_{i=1}^{n}a_{i}^{2}
      -2t\sum_{i=1}^{n}a_{i}b_{i}
      +t^{2}\sum_{i=1}^{n}b_{i}^{2}\\
    &=\Bigl(\sum_{i=1}^{n}b_{i}^{2}\Bigr)t^{2}
      -2\Bigl(\sum_{i=1}^{n}a_{i}b_{i}\Bigr)t
      +\sum_{i=1}^{n}a_{i}^{2}.
\end{align*}
Thus $f(t)=At^{2}+Bt+C$ with
\[
A=\sum_{i=1}^{n}b_{i}^{2},\qquad
B=-2\sum_{i=1}^{n}a_{i}b_{i},\qquad
C=\sum_{i=1}^{n}a_{i}^{2}.
\]
Note that $A\ge 0$.

\item \textbf{Non‑negativity of a quadratic.}  
A real quadratic $At^{2}+Bt+C$ with $A>0$ satisfies $At^{2}+Bt+C\ge0$ for all $t$ iff its discriminant is non‑positive:
\[
B^{2}-4AC\le 0 .
\]
(If $A=0$ the polynomial is linear; the degenerate cases are handled in step~6.)

\item \textbf{Discriminant.}  
Substituting $A,B,C$,
\begin{align*}
B^{2}-4AC
&=\bigl(-2\sum a_{i}b_{i}\bigr)^{2}
  -4\Bigl(\sum b_{i}^{2}\Bigr)\Bigl(\sum a_{i}^{2}\Bigr)\\
&=4\Bigl[\bigl(\sum a_{i}b_{i}\bigr)^{2}
         -\bigl(\sum a_{i}^{2}\bigr)\bigl(\sum b_{i}^{2}\bigr)\Bigr].
\end{align*}

\item \textbf{Deriving the inequality.}  
Since $f(t)\ge0$ for all $t$, the discriminant must satisfy $B^{2}-4AC\le0$. Dividing the expression in step~4 by $4$ yields
\[
\bigl(\sum_{i=1}^{n} a_{i}b_{i}\bigr)^{2}
\le
\bigl(\sum_{i=1}^{n} a_{i}^{2}\bigr)
\bigl(\sum_{i=1}^{n} b_{i}^{2}\bigr),
\]
which is the Cauchy–Schwarz inequality.

\item \textbf{Degenerate cases.}
\begin{itemize}
\item If $\displaystyle\sum_{i=1}^{n}a_{i}^{2}=0$, then $a_{i}=0$ for all $i$ and both sides of the inequality are $0$.
\item If $\displaystyle\sum_{i=1}^{n}b_{i}^{2}=0$, then $b_{i}=0$ for all $i$ and again the inequality reduces to $0\le0$.
\end{itemize}
Thus the inequality holds in all possible configurations.

\item \textbf{Equality condition.}  
Equality occurs precisely when $B^{2}-4AC=0$, i.e. when the quadratic $f(t)$ has a double root. For $A>0$ this happens iff there exists $\lambda\in\mathbb R$ such that
\[
a_{i}-\lambda b_{i}=0\quad\text{for all }i,
\]
or equivalently $a_{i}=\lambda b_{i}$ for every $i$. Indeed, solving $f'(\lambda)=0$ gives
\[
\lambda=\frac{\sum a_{i}b_{i}}{\sum b_{i}^{2}},
\]
and $f(\lambda)=0$ forces $a_{i}=\lambda b_{i}$ for each index. Conversely, if $a_{i}=\lambda b_{i}$ for all $i$, then each summand $(a_{i}-\lambda b_{i})^{2}$ is zero, so $f(\lambda)=0$ and the discriminant vanishes, yielding equality in the Cauchy–Schwarz inequality.

\item \textbf{Conclusion.}  
For arbitrary real sequences $\{a_i\}_{i=1}^{n}$ and $\{b_i\}_{i=1}^{n}$ the inequality
\[
\left(\sum_{i=1}^{n} a_{i}b_{i}\right)^{2}
\le
\left(\sum_{i=1}^{n} a_{i}^{2}\right)
\left(\sum_{i=1}^{n} b_{i}^{2}\right)
\]
holds, with equality exactly when the two vectors are linearly dependent (or one of them is the zero vector). \(\square\)
\end{enumerate}
\end{proof}